% JEDYNY PLIK, KTORY TRZEBA EDYTOWAC PRZED PRZYGOTOWANIEM PREZENTACJI.
% Jezyk: polish lub english. Typ pracy: engineering, master lub doctoral.
\newcommand{\presentationlanguage}{polish}
\newcommand{\thesistype}{engineering}

\newcommand{\universitypl}{Politechnika Opolska}
\newcommand{\universityen}{Opole University of Technology}
\newcommand{\facultypl}{Wydział Elektrotechniki, Automatyki i Informatyki}
\newcommand{\facultyen}{Faculty of Electrical Engineering, Automatic Control and Informatics}

\newcommand{\titlepl}{Tytuł pracy dyplomowej}
\newcommand{\titleen}{Diploma thesis title}
\newcommand{\authorname}{Imię i nazwisko}
\newcommand{\albumid}{000000}
\newcommand{\programmepl}{Nazwa kierunku}
\newcommand{\programmeen}{Study programme}
\newcommand{\supervisorname}{dr hab. inż. Imię Nazwisko, prof. uczelni}
\newcommand{\auxsupervisorname}{}
\newcommand{\city}{Opole}
\newcommand{\presentationyear}{2026}

% Logo KreativEU mozna wylaczyc, ustawiajac false.
\newcommand{\showkreativeulogo}{true}

% Ostatni slajd: title (powtorzenie strony tytulowej), thanks albo none.
\newcommand{\lastslidetype}{title}

% Pliki identyfikacji wizualnej.
\newcommand{\universitylogopl}{assets/branding/logo4-Politechnika_Opolska.pdf}
\newcommand{\universitylogoen}{assets/branding/logo-en-cropped.pdf}
\newcommand{\facultylogo}{assets/branding/Logo-WEAIL.png}
\newcommand{\kreativeulogo}{assets/branding/logos-KreativEU-vector-color.png}
